%% =====================================================================
%% 这是 main.tex —— 南昌大学本科生毕业设计（论文）模板【统一版 v1.2.4】
%% 编译：xelatex -> biber -> xelatex -> xelatex
%%       （v1.2.0 起参考文献改用 biblatex + biber，不再用 BibTeX）
%% ---------------------------------------------------------------------
%% 页面方案由类选项选择，同一套 ncubachelor.cls：
%%
%% · 默认（无选项）—— 单面打印版：
%%     页眉为当前章标题，居中，页眉线为文武线；章首页（含目录、
%%     摘要、声明、参考文献、附录等）一律 plain：无页眉，页码居中页脚。
%% · official（本示例所用）—— 官方样张版式：
%%     复刻 v1.1.5 锁定版：页眉章名居中 + 0.5pt 单横线，页码居中于
%%     页脚；各部分首页同样显示页眉（官方样张“全程有页眉”）。
%%     official 只决定页眉方案，可与 twoside / openright 叠加，例如
%%     双面印刷的规范组合（v1.2.2 起支持）：
%%       [official,twoside,openright]
%%     —— 双面装订、各构成要素奇数开页，但章首页仍 fancy（有页眉），
%%        不用 plain；空白填充页仍无页眉页脚。
%% · ncuhead —— 双面装订版（隐含 twoside 与 openright）：
%%     · 奇数页（右页）内侧：当前章标题
%%     · 偶数页（左页）内侧：南昌大学本科生毕业论文
%%     · 页码位于外侧：偶数页 LE、奇数页 RO
%%     · 页眉线为文武线（上粗下细双线）
%%     · 章首页及各部分首页一律 plain：无页眉，页码居中于页脚
%%     · 双面装订偏移默认 5mm（bindingoffset 可改，0mm 关闭）
%%     · 论文的全部构成要素一律“视为章”，从纸张正面（奇数物理页）
%%       开始，必要时自动补一张不编页码的空白页（见下）
%% ---------------------------------------------------------------------
%% 构成要素奇数开页（openright 开启时生效：ncuhead 隐含开启；
%% official,twoside,openright 组合同样生效）：
%%   声明、中文摘要、英文摘要、目录、正文各章、结论、参考文献、
%%   附录、致谢 —— 每一个都另起纸张正面；若上一要素恰好结束在纸张
%%   正面（奇数物理页），则在其后补一张不显示页码、页眉页脚的空白页
%%   （空白页仍计入页数）。因此封面之后会有若干空白页，例如：
%%     1 封面（正）  2 空白（背）
%%     3 声明（正）  4 空白（背）
%%     5 中文摘要（正） 6 空白（背）
%%     7 英文摘要（正） 8 空白（背）
%%     9 目录（正） 10 空白（背）
%%    11 正文第 1 章 …
%%   若希望要素紧接上一页、不补空白页，加 openany 关闭（见下）。
%% =====================================================================
\documentclass[official]{ncubachelor}
%% 参考文献数据库（v1.2.0 起用 biblatex 的 \addbibresource，
%% 取代旧的 \bibliography；文献列表在正文用 \printbibliography 输出）
%% 注意：文献后端是 biber（不是 bibtex）——编译链见 README 或 build.bat。
%% 若正文误写成旧写法 \bibliography{references}，v1.2.3 起会被接管并给
%% 中文提示（它不会像以前那样让文献表整个消失）。
\addbibresource{references.bib}
%% 常用类选项（可叠加，方括号内逗号分隔；official / ncuhead 为页眉
%% 方案、二选一后写者优先；twoside / openright 只管纸张，可任意叠加）：
%%   official       官方样张版式（本示例所用）：复刻 v1.1.5——页眉章名
%%                  居中 + 0.5pt 单横线，页码居中页脚，各部分首页同样
%%                  显示页眉；可与 twoside/openright 叠加（见上）
%%   ncuhead        另用版双面页眉（见上），隐含 twoside 与 openright
%%   twoside        双面打印（奇偶页眉在外侧显示章名，复旦思路），
%%                  并启用对称的内外侧边距与“空偶数页不显示页眉”
%%   openright      各构成要素从纸张正面（奇数页）开始（ncuhead 已隐式开启）
%%   openany        取消奇数开页：要素紧接上一页，不补空白页
%%   bindingoffset= 装订偏移（twoside 下默认 5mm，偶数页自动交替到内侧）
%%   blindreview    盲审版：封面/摘要中的姓名、学号、导师自动替换为 ***，
%%                  并清除 PDF 元数据中的作者信息（可与其他选项叠加）
%%   fontset=...    字体方案：auto(默认,自动探测)/windows/mac/founder/adobe/fandol/none
%%   tocdots        目录显示点线（v1.1.5 起默认开启，tocdots=false 关闭）

%% ---------------------------------------------------------------------
%% 一、论文元信息（按实际填写）
%% ---------------------------------------------------------------------
\ncusetup{
  info = {
    secret-level     = {公开},
    title            = {南昌大学本科生毕业论文模板使用演示},
    title*           = {Demonstration of the Nanchang University\\Undergraduate Thesis Template},
    college          = {物理},
    department       = {力学},
    major-class      = {物理学 2023 班},
    major            = {物理学},
    author           = {张三},
    studentid        = {402230220001},
    supervisor       = {李四},
    supervisor-title = {教授},
    start-date       = {23},
    stop-date        = {27},
    start-stop-date  = {2023 年 9 月—2027 年 6 月}
  }
}

%% ---------------------------------------------------------------------
%% 二、模板细节配置（可选，默认值如下）
%% ---------------------------------------------------------------------
%% 盲审细节（配合 blindreview 类选项）：
%%   blind = {
%%     text = {***},                  %% 遮蔽占位文本
%%     hide-acknowledgement = true,   %% 同时省略致谢页
%%   }
%% 定理环境样式：
%%   theorem = {
%%     header-font = \songti\bfseries,  %% 头部字体
%%     within      = chapter,           %% 编号锚点：chapter / section / none
%%   }
%%   proof = { qed = $\blacksquare$ }  %% 证明结束符

%% ---------------------------------------------------------------------
%% 三、其他宏包与自定义设置（按需添加）
%% ---------------------------------------------------------------------
% \usepackage{siunitx}
% \usepackage{algorithm,algorithmicx,algpseudocode}
% \usepackage{listings}

\usepackage{lipsum,zhlipsum}

%% ---------------------------------------------------------------------
%% 三、文档开始
%% ---------------------------------------------------------------------
\begin{document}

%% 封面（无页眉页脚）
\maketitle

%% 原创性申明与授权书（无页眉页脚）
\makedecaut

%% 前置部分：页码从中文摘要起用罗马数字 I, II, III …
\frontmatter

%% 中文摘要
\begin{abstract}
这是中文摘要。摘要应扼要说明研究目的、方法、结果和结论，
内容为小四号宋体、两端对齐、行距 1.35 倍、行首缩进 2 个汉字字符。

如果摘要较长，可以分段书写，也可以写到第二页。

本文得到了国家自然科学基金项目的资助。
\keywords{南昌大学；毕业论文；\LaTeX{} 模板}
\end{abstract}

%% 英文摘要
\begin{abstract*}
This is the English abstract. It should state briefly the purpose,
methods, results and conclusions of the research.

If the abstract is long enough, it may occupy more than one page.

This work was supported by the National Natural Science Foundation
of China.
\keywords{Nanchang University; bachelor thesis; \LaTeX{} template}
\end{abstract*}

%% 目录
\tableofcontents

%% ---------------------------------------------------------------------
%% 正文：页码从第一章起改用阿拉伯数字 1, 2, 3 …
%% ---------------------------------------------------------------------
\mainmatter

\chapter{引言}\label{chap:intro}

%% 若论文源于科研项目，可在正文第一页用无编号脚注标注项目来源
\blfootnote{本论文来源于科研项目“基于 \LaTeX{} 的本科毕业设计”，项目编号 1124。}

\section{研究背景与意义}\label{sec:background}

我们知道 lshort\cite{lshortcn} 是入门 \LaTeX{} 的必读文献。
可以花上几个小时阅读一下，能够避免被一些简单的问题耽误时间。
除此之外，还推荐刘海洋的《\LaTeX{} 入门》。

在社交媒体的强覆盖下，新闻信息的传播渠道也悄然发生了变化。
根据美国皮尤研究中心的调查结果\cite{pew}，社交媒体正在成为网络上
热点事件生成并发酵的源头，国内较为知名的新浪微博\footnote{本脚注
为小五号宋体，仅作排版示例。}即是典型代表。

\section{本文工作}

本文的主要工作如下：
\begin{enumerate}
  \item 依据《南昌大学本科生毕业设计（论文）书写式样》实现了
        封面、申明与授权书、中英文摘要、目录、正文、参考文献与致谢
        的自动排版；
  \item 提供了三线表、定理类环境、脚注等常用排版要素；
  \item 通过示例演示了模板的用法。
\end{enumerate}

\zhlipsum[3-7]

%% ---------------------------------------------------------------------
\chapter{图表与公式}\label{chap:figtab}

\section{表格}

表格标题置于表的上方，五号宋体加粗居中，表序与表名之间空一个汉字宽度；
表格内容为五号宋体，行距 1.35。表~\ref{tab:demo} 给出了一个三线表示例。

\begin{table}[htb]
  \centering
  \caption{用不同技术得到的带隙温度系数、$E_{g0}$、$\alpha$ 和 $T_0$ 的值}
  \label{tab:demo}
  \begin{tabular}{llccccc}
    \thickhline
    样品类型 & 实验方法 & $\mathrm{d}E_g/\mathrm{d}T$ (eV/K) &
    $E_{g0}$(eV) & $\alpha$ (eV/K) & $T_0$(K) & 文献 \\
    \hline
    GaN/Al$_2$O$_3$ & 光致发光 & $-5.32\times 10^{-4}$ & 3.503 &
    $5.08\times 10^{-4}$ & $-996$ & 61 \\
    GaN/Al$_2$O$_3$ & 光致发光 & —                      & 3.489 &
    $7.32\times 10^{-4}$ & 700   & 59 \\
    GaN/Al$_2$O$_3$ & 光吸收   & $-4.5\times 10^{-4}$   & $\sim$3.471 &
    $-9.3\times 10^{-4}$ & 772   & 63 \\
    \thickhline
  \end{tabular}
\end{table}

\section{插图}

图标题置于图的下方，五号宋体加粗居中，图序与图名之间空一个汉字宽度。
图~\ref{fig:demo} 给出了示例。请注意，浮动体不一定会留在插入位置，
因此不要使用“下图”“下表”这类表述，应使用“如图 1-1 所示”。

\begin{figure}[htb]
  \centering
  \includegraphics[width=.55\textwidth]{example-image}
  \caption{热风速计原理}
  \label{fig:demo}
\end{figure}

\section{公式}

对 Fourier 变换采用以下约定：
\begin{equation}
  (\mathcal{F}f)(\xi) = \hat{f}(\xi)
    = \frac{1}{\sqrt{2\pi}}\int_{-\infty}^{\infty}
      \mathrm{e}^{-ix\xi}f(x)\,\mathrm{d}x.
  \label{eq:fourier}
\end{equation}
根据这一标准定义，有
\[
  \lVert\hat{f}\rVert_{L^2} = \lVert f\rVert_{L^2},
  \qquad
  |\hat{f}(\xi)| \le (2\pi)^{-1/2}\lVert f\rVert_{L^1}.
\]
式~\eqref{eq:fourier} 中的积分顺序可以互换。

\section{定理类环境}

\begin{definition}[排版系统]
排版系统是把文字、公式、图表按规范组织成文档的软件系统。
\end{definition}

\begin{theorem}[可编译性]
只要元信息填写完整，本模板即可通过 XeLaTeX 顺利编译。
\end{theorem}

\begin{proof}
即得易见平凡，仿照上例显然。留作习题答案略，读者自证不难。
反之亦然同理，推论自然成立。略去过程 QED，由上可知证毕。
\end{proof}

\begin{lemma}
模板还提供引理、推论、命题、公理、假设、性质、注解等环境，
编号格式统一为“章.序”。
\end{lemma}

\begin{example}
在 \LaTeX{} 中书写“例”环境：本段文字即“例”环境的正文内容。
\end{example}

\begin{remark}
定理与证明的头部字体、编号锚点、证明结束符均可通过
\verb|\ncusetup{theorem=...}| 与 \verb|\ncusetup{proof=...}| 调整。
\end{remark}

\zhlipsum[3-7]

\section{双面页眉与页码（另用版演示）}

本节用于演示 \verb|ncuhead| 选项的双面页眉效果。本示例各章篇幅较短，
本节特意多写几段，使之后仍有余页，从而在纸张正反面各能看到一种完整的
页眉形态，便于对照检查。

偶数页（纸张背面，即左手页）的页眉内侧为“南昌大学本科生毕业论文”，
页码位于外侧左上角；奇数页（纸张正面，即右手页）的页眉内侧为当前章标题
（如“第二章　图表与公式”），页码位于外侧右上角；页眉线为文武线
（上粗下细的双线）。各部分的首页——章、目录、摘要、声明、参考文献、
附录——一律采用标准 plain 样式：不出现页眉，页码居中于页脚。

正文的页面设置、行距、首行缩进、图表与公式格式均与本模板单面版完全一致，
唯一的差别是页眉页脚与内外侧边距按纸张正反面自动交替。论文的全部构成要素
——声明、中文摘要、英文摘要、目录、正文各章、结论、参考文献、附录、致谢
——在本模板中一律“视为章”，都从纸张正面（奇数物理页）开始；若上一要素
恰好结束在纸张正面，则自动补一张不显示页码与页眉页脚的空白页（空白页仍
计入页数）。如需让要素紧接上一页、不补空白页，在类选项中加 \verb|openany|。

（以下段落为版面填充文字，可整段删除。）

页码的奇偶与纸张正反面始终保持一致：双面模式下，页码重新编号前会先做
纸张奇偶校准——若当前页落在纸张背面，则自动补一张空白页，使重新编号后
的首页落在纸张正面。否则内外侧边距、装订偏移与页眉的奇偶会与实际纸张
正反面相反。

装订偏移默认 5\,mm（A4 胶装通用值），偶数页自动交替到纸张内侧。若需要
调整，可使用 \verb|bindingoffset=8mm| 之类的选项；设为 \verb|0mm| 则关闭。
装订偏移消耗的宽度从版心中扣除，因此版心宽度、页眉线长度都会相应变化。

页眉文字与页码均为五号宋体（数字使用 Times New Roman），与官方《书写式样》
“页眉：中文宋体，五号”的要求一致。页码在正文起始处重新从 1 开始编号，
前置部分（中英文摘要、目录）使用罗马数字。

本页之后的页面即可看到奇数页页眉（本章标题在左、页码在右），再下一页
则为偶数页页眉（页码在左、“南昌大学本科生毕业论文”在右）。

若某处需要取消页眉（例如插图页、整页表格），可在该处使用
\verb|\thispagestyle{plain}|；若需要空白页不显示页眉页脚，可使用
\verb|\thispagestyle{empty}|。

模板的页眉内容取自 \verb|\leftmark|，因此章标题、无编号章标题（结论、
致谢）、参考文献标题都会自动进入页眉，无需手工设置。

（演示文字至此结束。）

%% ---------------------------------------------------------------------
\unnumberedchapter{结论}

\zhlipsum[3-7]

结论一般不参与章节编号，但要进目录并出现在页眉中，
因此这里使用无编号章命令 \verb|\unnumberedchapter{结论}|。

本章总结全文的主要工作与不足，并对后续研究进行展望。

%% ---------------------------------------------------------------------
\appendix
\chapter{附加材料}\label{chap:appendix}

附录与正文的章节标题用法没有区别，仅编号方式不同（附录 A、附录 B）。

%% ---------------------------------------------------------------------
%% 后置部分
%% ---------------------------------------------------------------------
\backmatter

\zhlipsum[3-5]

%% 参考文献（biblatex：\addbibresource 在导言区指定数据库）
\printbibliography

%% 致谢
\begin{acknowledgements}
\zhlipsum[3,5]
\signoff{张三}{2027 年 6 月}
\end{acknowledgements}

\end{document}
%% 结束
